\scanpage{218}
\begin{center}1 -- 47 --\end{center}\hfill n$^\circ$ 413

16. \uline{Mesure de Haar d'une variété de groupe}.

Conservons les notations de la fin du n$^\circ$ précédent.
L'espace $\Omega_e^n(G)$ est alors de dimension 1 et il en résulte que \uline{l'espace des formes différentielles invariantes à gauche et de degré} $n=\dim G$ \uline{est de dimension} 1.
On en obtient une \uline{base} en prenant une base $\omega_1,\ldots,\omega_n$ de $\Omega_e^1(G)$ et en prenant le produit extérieur des formes $I_{\omega_k}^g$.
Avec les notations de (30), on obtient en prenant $\omega_k=dx_k$, l'expression d'une forme invariante à gauche $I_\omega^g$ de degré maximum au voisinage de $e$:
\[
\tag{31}I_\omega^g(x)=(\operatorname{dét}\partial m_k/\partial y_j(x^{-1};x))\,dx_1\cdots dx_n.
\]
D'autre part, il est clair que l'espace des formes différentielles invariantes à gauche est invariant par translations à droite.
Il existe donc une fonction de classe $C^{r-1}$ sur $G$, soit $\delta$, telle que:
\[
\tag{32}I_\omega^g(xy)=\delta(y)I_\omega^g\qquad(x,y\in G,\ \omega\in\Omega_e^n(G)).
\]
Mais on a $\delta(y)I_\omega^g=I_\omega^g(xy)=I_\omega^g(y^{-1}xy)=I_{\omega'}^g$ avec $\omega'=(\operatorname{ad}y^{-1})_*(\omega)$.
D'où:
\[
\tag{33}\delta(y)=(\operatorname{dét}\operatorname{ad}y)^{-1}.
\]

Supposons maintenant que $K$ soit \uline{localement compact} et soit $\omega\in\Omega_e^n(G)$.
Il est clair que la mesure $\mu=|I_\omega^g|$ est une \uline{mesure de Haar à gauche} sur le groupe localement compact $G$.
Avec les notations de (30) et (31), on a donc pour $x$ voisin de $e$:
\[
\tag{34}d\mu(x)=\operatorname{mod}(\operatorname{dét}\partial m_k/\partial y_j(x^{-1};x))\,dx_1\cdots dx_n
\]
où mod est le \uline{module} du corps localement compact $K$ (Int. VII, \S\ 1, n$^\circ$ 10 et F. R. V.).
D'autre part, (32) et (33) montrent que le \uline{module} du groupe localement compact $G$ (Int. VII, \S\ 1, n$^\circ$ 3, déf. 3) n'est autre que
\[
\tag{35}\Delta(y)=\operatorname{mod}(\delta(y))=\operatorname{mod}(\operatorname{dét}\operatorname{ad}y)^{-1}.
\]

\margin{217}
